%! Trigonometric Equations and Inequalities — Unit 7, Lesson 7.3 — Homework (Answer Key)
\documentclass[10pt]{article}
\usepackage{precalculus-article}
\usepackage{precalculus-key}
\newcommand{\TallMath}[1]{$\displaystyle #1\rule[-1.4em]{0pt}{3.2em}$}

\begin{document}

\pageheader{Unit 7, Lesson 7.3}{Homework --- Answer Key}
\namedateperiod

\vspace{0.10in}

\begin{notesbox}{Problems 1--6}

\textbf{1.}\quad Solve each equation for all solutions (general solution).
Express answers in exact form using $\pi$.

\vspace{4pt}
(a)\quad $\sin\theta = \dfrac{\sqrt{2}}{2}$

\vspace{0.12in}
$\theta = $ \ans{$\dfrac{\pi}{4} + 2\pi k$ \quad and \quad $\dfrac{3\pi}{4} + 2\pi k$, for integer $k$}

\vspace{8pt}
(b)\quad $\cos\theta = -\dfrac{1}{2}$

\vspace{0.12in}
$\theta = $ \ans{$\dfrac{2\pi}{3} + 2\pi k$ \quad and \quad $\dfrac{4\pi}{3} + 2\pi k$, for integer $k$}

\vspace{8pt}
(c)\quad $\tan\theta = -1$

\vspace{0.12in}
$\theta = $ \ans{$-\dfrac{\pi}{4} + \pi k$, for integer $k$}

\vspace{0.14in}
\textbf{2.}\quad Solve $2\sin\theta - 1 = 0$ for all $\theta$ in $[0, 2\pi]$.
Show each step.

\vspace{4pt}
Step 1 --- Isolate $\sin\theta$: \ansline{$\sin\theta = \dfrac{1}{2}$}

\vspace{4pt}
Step 2 --- Reference angle: \ansline{$\arcsin\!\left(\dfrac{1}{2}\right) = \dfrac{\pi}{6}$}

\vspace{4pt}
Step 3 --- Identify quadrants and write solutions:
\ans{Sine is positive in Q1 and Q2. $\theta = \dfrac{\pi}{6}$ and $\theta = \pi - \dfrac{\pi}{6} = \dfrac{5\pi}{6}$.}

\vspace{0.14in}
\textbf{3.}\quad Solve $2\cos^2\theta - \cos\theta - 1 = 0$ for all $\theta$
in $[0, 2\pi]$.

\vspace{4pt}
Let $u = \cos\theta$. Factor: \ansline{$(2u + 1)(u - 1) = 0$}

\vspace{4pt}
Case 1: $u = $ \ans{$-\dfrac{1}{2}$} $\Rightarrow$ $\theta = $ \ans{$\dfrac{2\pi}{3}$ and $\dfrac{4\pi}{3}$}

\vspace{4pt}
Case 2: $u = $ \ans{$1$} $\Rightarrow$ $\theta = $ \ans{$0$ (and $2\pi$, same point)}

\begin{teachernote}
Full solution set in $[0,2\pi]$: $\theta \in \{0,\,\tfrac{2\pi}{3},\,\tfrac{4\pi}{3}\}$.
Students often miss $\theta=0$ because $\arccos(1)=0$ is on a quadrantal angle.
\end{teachernote}

\vspace{0.14in}
\textbf{4.}\quad $h(t) = 8\sin\!\left(\dfrac{\pi}{15}\,t\right) + 10$, $0 \le t \le 60$.

\vspace{4pt}
(a)\quad For what values of $t$ in $[0, 60]$ is $h(t) = 14$?

\ans{Set $8\sin\!\bigl(\tfrac{\pi}{15}t\bigr)+10=14$. Then $\sin\!\bigl(\tfrac{\pi}{15}t\bigr)=\tfrac{1}{2}$. Reference angle: $\tfrac{\pi}{6}$. In $[0,2\pi]$: $\tfrac{\pi}{15}t=\tfrac{\pi}{6}$ $\Rightarrow$ $t=2.5$\,s; and $\tfrac{\pi}{15}t=\tfrac{5\pi}{6}$ $\Rightarrow$ $t=12.5$\,s. Adding one period ($T=30$): $t=32.5$ and $t=42.5$\,s. All four are in $[0,60]$.}

\vspace{4pt}
(b)\quad For what values of $t$ in $[0, 60]$ is $h(t) \ge 14$?

\ans{$[2.5,\,12.5]\cup[32.5,\,42.5]$}

\vspace{0.14in}
\textbf{5. (AP-Style MC)}

\ans{(B)} \quad $\theta = \dfrac{\pi}{3}$ and $\theta = \dfrac{5\pi}{3}$
\textcolor{keyred}{\textbf{$\leftarrow$ correct}}

\begin{itemize}[leftmargin=1.4em, itemsep=2pt]
  \item (A) misses the Q4 solution.
  \item (C) gives the general form, not just $[0,2\pi]$.
  \item (D) and (E) are the solutions for $\sin\theta=1/2$ and $\cos\theta=-1/2$, respectively.
\end{itemize}

\vspace{0.14in}
\textbf{6. (AP-Style MC)}

\ans{(E)} \quad Both (A) and (B) \textcolor{keyred}{\textbf{$\leftarrow$ correct}}

\begin{itemize}[leftmargin=1.4em, itemsep=2pt]
  \item Both $\sin\theta=0.6$ and $\cos\theta=0.6$ have exactly two solutions in $(0,2\pi)$ (one in Q1, one in Q2 for sine; one in Q1, one in Q4 for cosine).
  \item Tangent is one-to-one per period so has only one solution per period.
\end{itemize}

\end{notesbox}

\vspace{0.08in}

\begin{extensionbox}
\textbf{Extension (Optional):} $2\sin^2\theta - \sin\theta - 1 = 0$

\vspace{4pt}
(a)\quad \ans{$(2\sin\theta + 1)(\sin\theta - 1) = 0$}

\vspace{6pt}
(b)\quad Case 1: $\sin\theta = -\tfrac{1}{2}$: \ans{$\theta = \dfrac{7\pi}{6} + 2\pi k$ and $\theta = \dfrac{11\pi}{6} + 2\pi k$}\\[4pt]
Case 2: $\sin\theta = 1$: \ans{$\theta = \dfrac{\pi}{2} + 2\pi k$}

\vspace{6pt}
(c)\quad In $[0, 4\pi]$: \ans{$\dfrac{\pi}{2},\,\dfrac{7\pi}{6},\,\dfrac{11\pi}{6},\,\dfrac{5\pi}{2},\,\dfrac{19\pi}{6},\,\dfrac{23\pi}{6}$ --- six solutions total.}
\end{extensionbox}

\vspace{0.08in}

\begin{tcolorbox}[colback=navylight, colframe=navy, arc=2mm,
    left=3mm, right=3mm, top=2mm, bottom=2mm]
\small\textbf{Preview --- Lesson 3.11: Secant, Cosecant, and Cotangent}\\
You have now solved equations involving sine, cosine, and tangent. In
Lesson~3.11 we meet three reciprocal functions --- secant, cosecant, and
cotangent --- and explore how their graphs and asymptotic behavior arise
directly from the functions you already know.
\end{tcolorbox}

\end{document}
